\begin{abstract}
We prove that strict comparison in the extended-functional sense implies
Jiang--Su absorption for separable simple nuclear non-elementary
$C^*$-algebras. This resolves the corresponding implication of the
Toms--Winter regularity problem, including nonunital algebras and allowing
unbounded traces. More generally, every separable nuclear $C^*$-algebra
whose Cuntz semigroup is almost unperforated and fully almost divisible
absorbs the Jiang--Su algebra.
\end{abstract}

\section{Introduction}
\label{sec:introduction}

Comparison asks whether numerical size determines order among positive
operators. Divisibility asks whether those operators can be split into
pieces of nearly equal size. For nuclear $C^*$-algebras, these questions
are tied to tensorial absorption of the Jiang--Su algebra $\Z$:
\[
 A\cong A\otimes_{\min}\Z.
\]
Jiang and Su constructed $\Z$ as a simple infinite-dimensional unital
algebra with the Elliott invariant of the complex numbers
\cite{JS1999}. Its absorption is a central regularity property in the
classification program. The purpose of this paper is to pass from
comparison and divisibility in the Cuntz semigroup to this tensorial
regularity, including algebras without a unit and algebras with ideals.

\subsection{Comparison and the simple theorem}

All algebras in this paper are complex. Let $\K$ be the compact
operators on a separable infinite-dimensional Hilbert space. For
$a,b\in(A\otimes\K)_+$, write $a\precsim b$ if
$v_j^*bv_j\to a$ in norm for some sequence $(v_j)$ in
$A\otimes\K$. Mutual comparison defines the Cuntz classes $[a]$.
Their ordered semigroup $\Cu(A)$ has addition given by orthogonal
sum and suprema of increasing sequences \cite{CEI2008}. A
\emph{functional} on $\Cu(A)$ is a map to $[0,\infty]$ preserving
zero, addition, order, and increasing suprema. Write
$\F(\Cu(A))$ for the cone of all these functionals.

We use the following precise extended-functional convention for
\emph{strict comparison}: for every $x,y\in\Cu(A)$,
\begin{equation}
\label{ext:eq:comparison}
\left.
\begin{aligned}
 \varphi(x)&\le\varphi(y)
       &&\text{for every }\varphi\in\F(\Cu(A)),\\
 \varphi(x)&<\varphi(y)
       &&\text{whenever }0<\varphi(y)<\infty
\end{aligned}
\right\}
\quad\Longrightarrow\quad x\le y.
\end{equation}
The second line concerns the same cone of functionals. Arithmetic at
zero and infinity is extended nonnegative arithmetic; in particular,
strict inequality is required precisely at positive finite target
values. An algebra is \emph{non-elementary} if it is not isomorphic to
the compact operators on any Hilbert space.

\begin{theorem}[Simple comparison and absorption]
\label{ext:thm:main}
Let $A$ be a separable simple nuclear non-elementary $C^*$-algebra.
If $A$ has strict comparison in the sense of
Equation~\eqref{ext:eq:comparison}, then $A\cong A\otimes_{\min}\Z$.
\end{theorem}

R\o rdam proved that $\Z$-absorption forces almost unperforation of
the finite-matrix Cuntz semigroup and, in the simple exact unital
setting, the corresponding trace comparison theorem
\cite[Theorem~4.5 and Corollary~4.6]{Rordam2004}. The reverse
implication is the strict-comparison-to-$\Z$-stability part of the
Toms--Winter regularity problem; see
\cite[Conjecture~5 and Problem~XVIII]{Problems2025}.
Theorem~\ref{ext:thm:main} gives a positive resolution of this
implication with the extended-functional convention above. Its
nonunital formulation includes stably projectionless algebras and
allows unbounded traces.

The comparison argument has a stronger intermediate scope.
Theorem~\ref{ext:thm:divisibility} proves, for every simple exact
non-elementary $A$ satisfying Equation~\eqref{ext:eq:comparison},
that each $x\in\Cu(A)$ and integer $N\ge1$ admit a single class $u$
with $Nu\le x\le(N+1)u$. Neither separability nor nuclearity enters
that theorem. Nuclearity is used in constructing central maps from
this order-theoretic information.

\subsection{The nonsimple absorption mechanism}

The first order property we require is \emph{almost unperforation}:
\begin{equation}
 (l+1)x\le ly\quad\Longrightarrow\quad x\le y
 \qquad(x,y\in\Cu(A),\ l\ge1).
 \label{eq:pre:au}
\end{equation}
The second is \emph{full almost divisibility}: for every
$x\in\Cu(A)$ and integer $N\ge1$, there is $u\in\Cu(A)$ such that
\begin{equation}
 Nu\le x\le(N+1)u.
 \label{eq:pre:div}
\end{equation}
Here $l$ is also an integer. We call the conjunction of these two
properties \textup{(C)}. The same class $u$ must satisfy both
inequalities in Equation~\eqref{eq:pre:div}.

The comparison convention in Equation~\eqref{ext:eq:comparison} is
equivalent to almost unperforation of $\Cu(A)$
\cite[Proposition~6.2]{ERS2011}. Lemma~\ref{lem:comparison-convention}
explains this equivalence, including the zero and infinite values.
Thus the extended-functional formulation agrees with the usual
Cuntz-semigroup comparison condition.

To state the central conclusion, fix a free ultrafilter $\omega$ on
$\mathbb N$. The norm ultrapower and its central annihilator quotient
are
\[
 A_\omega=\ell^\infty(\mathbb N,A)/
       \{(a_n):\lim_{n\to\omega}\|a_n\|=0\},\qquad
 F_\omega(A)=(A_\omega\cap A')/\Ann(A),
\]
where $A$ is the constant copy and
$\Ann(A)=\{z\in A_\omega:zA=Az=0\}$. For separable $A$,
a sequential approximate unit represents the unit of $F_\omega(A)$.
The zero algebra absorbs $\Z$ immediately; the central-map statement
below is formulated for nonzero algebras.

\begin{theorem}[Absorption from condition \textup{(C)}]
\label{thm:condition-C}
Let $A$ be a nonzero separable nuclear $C^*$-algebra satisfying
\textup{(C)}. For every fixed free ultrafilter $\omega$ on
$\mathbb N$, there is a unital $*$-homomorphism
$\Z\to F_\omega(A)$. Consequently $A\cong A\otimes_{\min}\Z$.
\end{theorem}

\begin{corollary}[The canonical Cuntz-semigroup map]
\label{cor:canonical}
Let $A$ be a separable nuclear $C^*$-algebra, and let
$\iota:A\to A\otimes_{\min}\Z$ be $\iota(a)=a\otimes1_\Z$.
If $\Cu(\iota):\Cu(A)\to\Cu(A\otimes_{\min}\Z)$ is an
isomorphism, then $A\cong A\otimes_{\min}\Z$.
\end{corollary}

The isomorphism in this hypothesis is the map induced by $\iota$.
Corollary~\ref{cor:canonical} resolves positively implication
(iii)$\Rightarrow$(ii) in the nonsimple Toms--Winter problem
\cite[Section~22, Problem~LXXV]{Problems2025}. It applies to every
separable nuclear $A$ with the indicated canonical map; an additional
restriction on its elementary ideal subquotients is unnecessary.

Condition \textup{(C)} depends only on the ordered-semigroup structure.
Consequently, the same absorption theorem also removes the requirement
that the Cuntz-semigroup isomorphism be the canonical one.

\begin{corollary}[An abstract Cuntz-semigroup criterion]
\label{cor:abstract-cu}
Let $A$ be a separable nuclear $C^*$-algebra. Then
\[
 A\cong A\otimes_{\min}\Z
 \quad\Longleftrightarrow\quad
 \Cu(A)\cong\Cu(A\otimes_{\min}\Z)
 \text{ in the category }\mathbf{Cu}.
\]
The isomorphism on the right is abstract: it is not required to be
induced by the first-factor embedding, or by any $*$-homomorphism.
\end{corollary}

Here the category $\mathbf{Cu}$ retains zero, addition, order, increasing
sequential suprema, and the way-below relation of the Cuntz semigroup.
Corollary~\ref{cor:abstract-cu} answers positively the abstract
Cuntz-semigroup question in
\cite[Section~9, Problem~XXVI]{Problems2025}, without assuming
unitality, simplicity, or stable finiteness. Separability and nuclearity
remain hypotheses.

\subsection{History and the main obstacles}

The comparison of positive elements and their dimension functions goes
back to Cuntz and Blackadar--Handelman
\cite{Cuntz1978,BlackadarHandelman1982}. The completed Cuntz semigroup
of Coward--Elliott--Ivanescu \cite{CEI2008} makes increasing suprema
part of the invariant. Elliott--Robert--Santiago's correspondence
\cite{ERS2011} identifies its functionals with lower semicontinuous
extended quasitraces. For exact algebras, Haagerup's additivity theorem
\cite{Haagerup2014}, in the extended form of Blanchard--Kirchberg
recalled in \cite[Remark~4.3]{ERS2011}, makes these quasitraces traces.
The associated rank is
$d_\tau(a)=\lim_n\tau(a^{1/n})$: the functional corresponding to
$\tau$ sends $[a]$ to $d_\tau(a)$.

The reverse comparison implication has a substantial history of
partial results. Matui and Sato proved it for unital separable simple
nuclear infinite-dimensional algebras with finitely many extremal
traces, using excision in small central sequences and property
\textup{(SI)} \cite[Theorem~1.1]{MatuiSato2012}.
Sato, Kirchberg--R\o rdam, and Toms--White--Winter extended the
tracial-boundary approach to nonempty compact extreme boundaries of
finite covering dimension \cite[Corollary~1.2]{Sato2012},
\cite[Corollary~7.9]{KirchbergRordam2014},
\cite[Corollary~4.7]{TWW2015}.
Castillejos, Evington, Tikuisis, and White isolated the remaining
uniform central divisibility in trace: for simple separable unital
nuclear algebras with nonempty tracial state space, strict comparison
together with uniform property~$\Gamma$ is equivalent to
$\Z$-stability \cite[Theorem~5.6]{CETW2022}. Uniform property~$\Gamma$
requires approximately central partitions that split the trace values
of finitely many elements evenly, uniformly over all traces.
The present argument constructs
central maps from condition \textup{(C)} without imposing a
tracial-boundary or uniform-$\Gamma$ hypothesis.

The Toms--Winter problem also links these properties to finite nuclear
dimension, a finite-colour form of completely positive approximation.
Winter's absorption theorem and the converse of Castillejos, Evington,
Tikuisis, White, and Winter establish the equivalence of finite nuclear
dimension and $\Z$-stability in the simple separable unital nuclear
non-elementary setting \cite[Corollary~6.3]{Winter2012}
\cite[Theorem~A]{CETWW2021}. The comparison implication is the
remaining direction addressed here.

The distinction between comparison and divisibility is already
visible in the range of possible rank functions of positive elements.
Thiel proved that strict comparison yields almost divisibility and
the canonical Cuntz-semigroup isomorphism for separable simple unital
non-elementary algebras of stable rank one
\cite[Corollary~8.12]{Thiel2020}. Theorem~\ref{ext:thm:divisibility}
uses exactness without a stable-rank-one hypothesis.
The fractional spectral discretization in Dadarlat--Toms
\cite[Theorem~5.1 and its proof]{DadarlatToms2010} follows earlier
rank-realization ideas of Brown--Perera--Toms \cite[Section~5]{BPT2008}.
Our first construction moves trace values
into a hereditary subalgebra of small rank before taking integer
spectral cuts. The resulting error is at most one support rank,
regardless of the norm of the transferred element or the number of
spectral levels. Proposition~\ref{prop:unital-approx} gives the unital
uniform-trace version of this elementary construction.

For nonunital algebras, there may be no normalized trace space on
which all errors can be bounded. The substitute is one fixed nonzero
spectral cut of the element being divided. Simplicity makes its rank
positive and finite for every functional that is finite and positive
somewhere, even when the original element has infinite rank.
We measure all rounding errors relative to this same cut and make them
summable. Zero functionals and functionals infinite on every nonzero
class are handled separately. This keeps the full extended trace cone
available without requiring a common normalization; the precise
three-way distinction is proved in Lemma~\ref{ext:lem:trichotomy}.

Comparison and divisibility are also studied together through the
notion of pureness. For the completed semigroup, the usual almost
divisibility condition asks for $Nu\le x$ and $x'\le(N+1)u$ whenever
$x'$ is compactly contained in $x$; together with almost unperforation
this defines pureness \cite{APT2018}. Full almost divisibility as
used here implies that condition.
For separable simple unital non-elementary algebras of locally
finite nuclear dimension, Winter proved that pureness implies
Jiang--Su absorption \cite[Corollary~6.2]{Winter2012}. Under the same
local nuclear-dimension hypothesis, he also obtained an abstract
criterion using $W(A)$, the Cuntz semigroup of positive elements in
finite matrices over $A$ \cite[Corollary~6.4]{Winter2012}.
Corollary~\ref{cor:abstract-cu} instead uses the completed invariant
$\Cu(A)$ and has no local nuclear-dimension assumption.
The tensor-product approach to regularity of abstract Cuntz semigroups
was developed by Antoine, Perera, and
Thiel \cite{APT2018}; more recent work of Antoine, Perera, Thiel,
and Vilalta characterizes pureness by controlled comparison and
functional divisibility \cite[Theorem~C]{APTV2024}.
Regularity of $\Cu(A)$ and regularity of the central
sequence algebra are distinct inputs. In particular,
Perera, Thiel, and Vilalta characterize Jiang--Su absorption by
pureness of the central sequence algebra for every separable $A$
\cite[Theorem~3.6]{PTV2025}. Theorem~\ref{thm:condition-C}
constructs the required central structure from the stated order
properties of $\Cu(A)$ itself.

For the central construction, a \emph{matrix cone} means a completely
positive contractive order-zero map from $M_k$: it preserves
orthogonality of positive elements. Its \emph{weight} is the image of
$1_{M_k}$. Norm-central cones take values in $A_\omega\cap A'$,
with coordinate maps whose centrality and order-zero errors tend to
zero along $\omega$.

Robert and Tikuisis developed a nonsimple route from regularity to
central structure using finite nuclear dimension
\cite[Theorem~1.1]{RobertTikuisis2017}. For separable algebras of
finite nuclear dimension, they characterized absorption by the
existence of two full orthogonal elements in the central sequence
algebra \cite[Theorem~1.2]{RobertTikuisis2017}.
Our hypothesis is instead the order condition \textup{(C)}, with no
bound on nuclear dimension.
The nonsimple problem has a different obstruction. A positive element
can be large in every representation of interest without a single
uniform integer controlling all the associated ideal comparisons.
At the same time, traces may be finite only after hereditary
localization. The proof treats these two forms of largeness
separately. Pure-state compression and support shrinking construct
norm-visible central matrix cones. Convex order-zero nuclear
approximations, originating in the work of Hirshberg, Kirchberg,
and White \cite{HKW2012} and refined by Brown, Carri\'on,
and White \cite{BCW2019}, supply cones large for bounded traces.
We localize the latter estimates and then remove a relative trace
ideal to obtain norm centrality.

The matrix-cone calculus uses Winter and Zacharias's structure
theorem for order-zero maps \cite{WZ2009}. The final packing step is
an instance of Gabe's approximate domination theorem
\cite{Gabe2020}; its weighted-inclusion form is proved here.
Choi--Effros lifting and Arveson extension allow us to retain
completely positive maps during diagonalization
\cite{CE1976,Arveson1969}. The absorption step uses the structural
properties of $\Z$ established by Jiang--Su, Toms--Winter, and
R\o rdam--Winter \cite{JS1999,TW2007,RW2010}. Nawata's nonunital
central dimension-drop criterion \cite[Proposition~5.1]{Nawata2013}
is implemented here by an explicit lifting and intertwining argument.

\subsection{Proof strategy and organization}

Figure~\ref{fig:proof-flow} separates the two routes into condition
\textup{(C)} from the common central-sequence argument. The full
proof of Theorem~\ref{thm:condition-C} is contained in this paper.
Its concrete target, for each $p\ge2$, is a matrix cone
$\alpha:M_p\to F_\omega(A)$ and an element $s$ with
\[
 s^*s=1-\alpha(1),\qquad \alpha(e_{11})s=s.
\]
Thus $s$ fills the cone's defect inside one diagonal corner.
These relations give a unital map from the prime dimension-drop
algebra $I(p,p+1)$, whose continuous functions on $[0,1]$ take values in
$M_p\otimes M_{p+1}$ and have endpoint values in
$M_p\otimes1$ and $1\otimes M_{p+1}$, respectively.
Such maps supply the central copy of $\Z$.

\begin{figure}[H]
\centering
\setlength{\tabcolsep}{7pt}
\begin{tabular}{c@{\qquad}c}
\fbox{\parbox{.39\linewidth}{\centering Simple exact non-elementary algebra:\\
extended-functional comparison}} &
\fbox{\parbox{.39\linewidth}{\centering Separable nuclear algebra:\\
abstract $\Cu(A)\cong\Cu(A\otimes\Z)$}}\\[5pt]
rounding and comparison arithmetic $\downarrow$ & $\downarrow$ transfer of tensor regularity\\[5pt]
\multicolumn{2}{c}{\fbox{\parbox{.67\linewidth}{\centering
Condition (C): almost unperforation and full almost divisibility}}}\\[7pt]
\multicolumn{2}{c}{\textit{Under separability and nuclearity:}}\\[5pt]
\fbox{\parbox{.39\linewidth}{\centering Norm-visible cones\\
fix the coordinate multiplicities}} &
\fbox{\parbox{.39\linewidth}{\centering Localized tracial cones\\
choose relative errors afterwards}}\\[5pt]
\multicolumn{2}{c}{$\downarrow$ bounded comparison witnesses and central slots}\\[5pt]
\multicolumn{2}{c}{$\downarrow$ central defect filler and dimension-drop maps}\\[5pt]
\multicolumn{2}{c}{$\Z\longrightarrow F_\omega(A)\quad\Longrightarrow\quad
 A\cong A\otimes\Z$}
\end{tabular}
\caption{The order of the constructions. The simple comparison route
supplies condition (C) already under exactness. Both absorption routes
then use the same separable nuclear construction. The norm-cone
multiplicities are fixed before the localized trace accuracies.}
\label{fig:proof-flow}
\end{figure}

Section~\ref{sec:preliminaries} introduces the common Cuntz-semigroup
and trace conventions, together with the support calculus for spectral
cuts. Sections~\ref{ext:sec:preliminaries}--\ref{ext:sec:divisibility}
then prove the simple exact almost-divisibility theorem. Trace values
are transferred to successively smaller hereditary supports, rounded
at integer levels, and summed in one norm-convergent stable block sum.
A fixed nonzero cut controls the summable errors even when the original
class has infinite functional value. Comparison supplies the two sides
of the divisibility sandwich, completing the simple route to
condition \textup{(C)}.

Section~\ref{sec:support-tools} begins the common absorption argument.
It derives condition \textup{(C)} from the canonical Cuntz-semigroup
map and develops its consequences for rank comparison and divisible
matrix cones. It also records the ultrapower and bounded
diagonalization tools used in the central constructions.
Sections~\ref{shr:section} and~\ref{norm:section} construct the
norm-visible cones. Two-support shrinking preserves the
representations detecting a prescribed positive cut while making
cross terms small. Pure-state word models can then be combined into
approximately central matrix cones. Their norm visibility supplies
finite coordinate ideal multiplicities.
Sections~\ref{sec:tracial} and~\ref{sec:localization} construct the
tracial cones. Two support compressions keep their errors independent
of the number of nuclear approximation summands; hereditary
localization and a relative trace ideal turn tracial centrality into
norm centrality.

Section~\ref{sec:comparison} combines these constructions in that
order. Its trace accuracies depend on the already fixed ideal
multiplicities. Spectral margins first produce exact coordinate
Cuntz inequalities and then bound the comparison witnesses
independently of matrix size. This permits a slowly growing supply
of orthogonal equivalent central slots.
Section~\ref{cen:section} packs finite simultaneous compressions
into those slots; approximate-unit tests make the packed element
bounded and central. Section~\ref{abs:section} turns the resulting
relations into dimension-drop maps, constructs a single unital
central copy of $\Z$, and performs the nonunital intertwining.
Section~\ref{sec:assembly} assembles the four stated results.
